对于磁偶极跃迁，本节与第2.7节统一采用磁矩型算符 $\mathcal M_{\mu}(M1)$，其核磁子为 SI 量 $\mu_N=e\hbar/(2m_p)$，并记相应的约化跃迁几率为 $B_{\mu}(M1)$，单位为 $\mu_N^2$。Suhonen 的电荷型磁算符满足 $\mathcal M_{S}(M1)=\mathcal M_{\mu}(M1)/c$，故 $B_{S}(M1)=B_{\mu}(M1)/c^2$。但按各自单位表示的无量纲数值相同：
\begin{equation}
\frac{B_{S}(M1)}{(\mu_N/c)^2}=\frac{B_{\mu}(M1)}{\mu_N^2}\equiv b_{M1}.
\end{equation}
原书 Eq.~(6.14) 对 $\lambda=1$ 给出~\cite{book-suhonen-from-nucleons-to-nucleus}
\begin{equation}\label{eq:5-9-cm1}
\lambda_\gamma(M1)=C_{M1}\left(\frac{E_\gamma}{\mathrm{MeV}}\right)^3 b_{M1},\qquad
C_{M1}=\frac{6.080\times10^{20}(2/9)}{197.33^3}\,\mathrm{s}^{-1}\simeq 1.7584\times10^{13}\,\mathrm{s}^{-1}.
\end{equation}
因此，原书采用 $(\mu_N/c)^2$ 单位时的 $b_{M1}$ 可直接作为常规 $\mu_N^2$ 单位下的数值使用。由光子分量反演磁矩型约化跃迁几率，得
\begin{equation}\label{eq:5-9-bm1-extracted}
B_{\mu}(M1)=\mu_N^2\frac{\lambda_\gamma(M1)}{C_{M1}(E_\gamma/\mathrm{MeV})^3}
=\mu_N^2\frac{b_k}{\tau(1+\alpha_k)(1+\delta^2)C_{M1}(E_\gamma/\mathrm{MeV})^3}.
\end{equation}
高斯单位制中核磁子定义为 $e_{\mathrm G}\hbar/(2m_pc)$，其中电荷 $e_{\mathrm G}$ 亦须采用高斯单位；不可将该定义与 SI 电荷及带量纲的光速数值直接混用。比较文献时应先确认算符和所用核磁子，再比较以相应单位归一的数值。
